\documentclass[11pt,a4paper]{article}

\usepackage[T1]{fontenc}
\usepackage[utf8]{inputenc}
\usepackage{lmodern}
\usepackage{microtype}
\usepackage[margin=2.3cm]{geometry}
\usepackage{enumitem}
\usepackage{xcolor}
\usepackage{xurl}
\usepackage{hyperref}

\hypersetup{
  colorlinks=true,
  linkcolor=blue!50!black,
  urlcolor=blue!50!black,
  pdftitle={LM-IDNet Code Quality and Structure Audit}
}
\setlength{\parindent}{0pt}
\setlength{\parskip}{0.45em}
\setlength{\emergencystretch}{2em}
\Urlmuskip=0mu plus 1mu
\setlist[enumerate]{leftmargin=1.8em,itemsep=1em,topsep=0.4em,
  label=\textbf{\arabic*.}}

\title{\textbf{LM-IDNet Code Quality and Structure Audit}}
\author{}
\date{19 October 2026}

\begin{document}
\maketitle

\textbf{What this repo does:} LM-IDNet turns device traffic into count windows,
learns a normal model, and produces explainable anomaly alerts with optional
guarded updates. The assumed load is one researcher running batch experiments
and one streaming process per device state directory.

The overall structure is understandable. Processing, modelling, runtime, and
saved state have clear roles, and batch and streaming share the detector.
\textbf{601 tests passed}, including container checks. The cases below were also
reproduced using temporary files. \textbf{No code was changed during the audit.}

\section*{Must fix}

\begin{enumerate}[series=findings]
\item \begin{minipage}[t]{\linewidth}
\textbf{A valid high-volume window stops the detector}

\path{src/lm_idnet/models/scoring_stage.py:70}

\textbf{What this is:} Likelihood scoring runs before the volume checks.

\textbf{Problem:} Five million TCP packets in a Chromecast window trigger
\path{NumericalPrecisionError}. The streaming process emits no event and
leaves the next window unprocessed.

\textbf{Fix:} Handle precision failures per window, report them explicitly,
exclude that window from learning, and continue. Preserve any independently
available volume evidence.

\textbf{If we skip it:} A traffic burst can interrupt monitoring.
\end{minipage}

\item \begin{minipage}[t]{\linewidth}
\textbf{Normal preprocessing accepts incomplete inputs}

\path{src/lm_idnet/processing/processing_manager.py:59}

\textbf{What this is:} The normal preprocessing command converts configured
captures into windows.

\textbf{Problem:} Missing captures produce a warning and are skipped. A capture
with a truncated record header also produced processed JSON, although the
separate validator correctly rejected it.

\textbf{Fix:} Enforce complete-file checks during preprocessing and fail on
missing configured captures, preserving support for both capture formats.

\textbf{If we skip it:} A successful command can leave incomplete or stale
processed data.
\end{minipage}

\item \begin{minipage}[t]{\linewidth}
\textbf{Batch scoring accepts incompatible calibration settings}

\path{src/lm_idnet/artifacts.py:148}

\textbf{What this is:} This loads the saved model and threshold.

\textbf{Problem:} Changing the configured quantile to \texttt{0.5} still accepts
a threshold calibrated at \texttt{0.01}. An increased minimum sample requirement
is also ignored.

\textbf{Fix:} Check the saved quantile, calibration captures, and sample count
against the configuration, as state restoration already does.

\textbf{If we skip it:} An experiment can run with different calibration settings
from those its configuration declares.
\end{minipage}

\item \begin{minipage}[t]{\linewidth}
\textbf{The ordinary evaluator accepts mismatched evidence}

\path{src/lm_idnet/evaluation/evaluation_stage.py:45}

\textbf{What this is:} The CLI matches prediction events to labels.

\textbf{Problem:} Events naming another device were accepted. Removing an attack
window's prediction also produced 100\% accuracy on the remaining matched window.

\textbf{Fix:} Check device identity and require coverage of observed windows.
Allow omissions only for explicitly missing source coverage.

\textbf{If we skip it:} Wrong or incomplete event files can produce misleading
summaries.
\end{minipage}

\item \begin{minipage}[t]{\linewidth}
\textbf{Failed writes can destroy saved research artifacts}

\path{src/lm_idnet/artifacts.py:164}

\textbf{What this is:} Ordinary model and threshold saves overwrite their
destination directly.

\textbf{Problem:} A simulated disk-full failure left a previously valid model
containing only \texttt{\{}. The protected detector-state writer already handles
this more safely.

\textbf{Fix:} Write a temporary file and replace the destination only after
successful completion.

\textbf{If we skip it:} Interrupted training or calibration can destroy the
previous usable artifact.
\end{minipage}

\item \begin{minipage}[t]{\linewidth}
\textbf{F1 is missing when every attack is missed}

\path{src/lm_idnet/evaluation/evaluation_stage.py:115}

\textbf{What this is:} F1 combines precision and recall.

\textbf{Problem:} With an attack present and no detections, the CLI returns
\texttt{null} instead of \texttt{0}. The experiment helper handles this case
correctly.

\textbf{Fix:} Calculate F1 directly from the confusion counts and test the
all-missed case.

\textbf{If we skip it:} The least successful runs can disappear from comparisons
as unavailable results.
\end{minipage}

\end{enumerate}

\section*{Should fix}

\begin{enumerate}[resume=findings]
\item \begin{minipage}[t]{\linewidth}
\textbf{Batch inputs lack the stream's timestamp safeguards}

\path{src/lm_idnet/models/modeling_stage.py:85}

\textbf{What this is:} Training loads count windows from processed files.

\textbf{Problem:} Repeating the same timestamped window twice is accepted as two
observations. Streaming already rejects duplicates and overlaps.

\textbf{Fix:} Add a shared check for ordered, non-overlapping windows within
each capture.

\textbf{If we skip it:} Duplicated records can distort fitting, calibration,
or adaptation timing.
\end{minipage}

\item \begin{minipage}[t]{\linewidth}
\textbf{Two production dependencies have no application callers}

\path{pyproject.toml:11}

\textbf{What this is:} Every installation includes \texttt{mpmath} and
\texttt{tabulate}.

\textbf{Problem:} Neither the installed application nor its tests imports
either package. The original research capsule declares its own
\texttt{mpmath} requirement separately.

\textbf{Fix:} Remove both from application dependencies; retain the capsule's
environment.

\textbf{If we skip it:} Installations carry unnecessary packages and version
constraints.
\end{minipage}

\item \begin{minipage}[t]{\linewidth}
\textbf{The production native file still contains unused demonstration code}

\path{src/lm_idnet/algorithms/native/LM_time_lib.c:23}

\textbf{What this is:} The compiled library includes old demonstration entry
points alongside the numerical implementation.

\textbf{Problem:} Functions such as \texttt{pippero()} read \path{gaucho.txt};
other unused functions load experimental datasets. No current application
caller uses these functions.

\textbf{Fix:} Remove these entry points from the production copy. Preserve the
original research capsule, numerical routines, and attribution.

\textbf{If we skip it:} Maintainers must distinguish obsolete experiments from
functioning detector code.
\end{minipage}

\end{enumerate}

\section*{Nice to have}

\begin{enumerate}[resume=findings]
\item \begin{minipage}[t]{\linewidth}
\textbf{Remove configuration fields that do nothing}

\path{src/lm_idnet/config.py:169}

\textbf{What this is:} \texttt{time\_col} and \texttt{protocol\_col} look like
preprocessing controls.

\textbf{Problem:} No packet reader uses them. They only change fingerprint
metadata.

\textbf{Fix:} Remove them from active configurations, schemas, and corresponding
fingerprint fields.

\textbf{If we skip it:} Users can change an apparent input setting without
changing processing.
\end{minipage}

\item \begin{minipage}[t]{\linewidth}
\textbf{Delete the unused distance-metric module}

\path{src/lm_idnet/evaluation/metrics.py:1}

\textbf{What this is:} It defines three small distance functions.

\textbf{Problem:} None has a caller in the application or tests; the research
script calculates its own distances.

\textbf{Fix:} Delete the unused module and update its documentation references.

\textbf{If we skip it:} Another implementation remains for readers to interpret
and maintain.
\end{minipage}

\item \begin{minipage}[t]{\linewidth}
\textbf{Remove a broken package export}

\path{src/lm_idnet/processing/__init__.py:4}

\textbf{What this is:} The package advertises an \texttt{aggregation} module.

\textbf{Problem:} That module does not exist.
\texttt{from lm\_idnet.processing import *} raises \texttt{AttributeError}.

\textbf{Fix:} Delete the stale export.

\textbf{If we skip it:} The advertised package interface remains broken.
\end{minipage}

\end{enumerate}

\bigskip
\begin{samepage}
\textbf{Verdict:} Keep the existing architecture. Fix the stream failure and
input/evaluation boundaries first, then remove the unused scaffolding.

\textbf{Lean:} Approximately 175 fewer lines and two fewer production
dependencies are possible.

\textbf{Not checked:} Formal mathematical proofs, every original
research-capsule script, exhaustive raw captures, physical hardware,
or a fresh image build.
\end{samepage}

\end{document}
